%Data institusi

%\input{Dozentdaten}

%\renewcommand{\fachbereich}{Fachbereich}

%\renewcommand{\dozent}{Prof. Dr. . }

%Data ujian

\renewcommand{\klausurgebiet}{ }

\renewcommand{\klausurtyp}{ }

\renewcommand{\klausurdatum}{ . 20}

\klausurvorspann {\fachbereich} {\klausurdatum} {\dozent} {\klausurgebiet} {\klausurtyp}

%Data untuk tabel nilai berikut

\renewcommand{\aeins}{  3  }

\renewcommand{\azwei}{  3   }

\renewcommand{\adrei}{  3  }

\renewcommand{\avier}{  0  }

\renewcommand{\afuenf}{  2  }

\renewcommand{\asechs}{  0  }

\renewcommand{\asieben}{  3  }

\renewcommand{\aacht}{  7  }

\renewcommand{\aneun}{  0  }

\renewcommand{\azehn}{  7  }

\renewcommand{\aelf}{  4  }

\renewcommand{\azwoelf}{  3  }

\renewcommand{\adreizehn}{  3   }

\renewcommand{\avierzehn}{  7  }

\renewcommand{\afuenfzehn}{  2  }

\renewcommand{\asechzehn}{  9  }

\renewcommand{\asiebzehn}{  7  }

\renewcommand{\aachtzehn}{  63  }

\renewcommand{\aneunzehn}{    }

\renewcommand{\azwanzig}{    }

\renewcommand{\aeinundzwanzig}{    }

\renewcommand{\azweiundzwanzig}{    }

\renewcommand{\adreiundzwanzig}{    }

\renewcommand{\avierundzwanzig}{    }

\renewcommand{\afuenfundzwanzig}{    }

\renewcommand{\asechsundzwanzig}{    }

\punktetabellesiebzehn

\klausurnote

\newpage

\setcounter{section}{0}

\inputaufgabegibtloesung
{3}
{

Definisikan istilah-istilah berikut
\zusatzklammer {yang dicetak miring} {} {}.
\aufzaehlungsechs{\stichwort {Pemetaan Gauss} {}
dari suatu hiperpermukaan diferensiabel

\mavergleichskette
{\vergleichskette
{ Y
}
{ \subseteq }{ \R^n
}
{  }{
}
{    }{
}
{   }{
}
}
{}{}{.}

}{Suatu
\stichwort {peta topologis} {}
pada suatu
\definitionsverweis {manifold topologis}{}{}
$M$.

}{Suatu $C^1$-\stichwort {manifold diferensiabel} {} $M$.

}{\stichwort {Pangkat eksterior} {} ke-$n$ dari suatu
$K$-\definitionsverweis {ruang vektor}{}{}
$V$
\zusatzklammer {cukup tuliskan simbolnya} {} {.}

}{Suatu
\stichwort {isometri lokal} {}
\maabb {\varphi} {L} {M
} {}
antara manifold-manifold Riemann.

}{Suatu koneksi
\stichwort {terintegralkan secara lokal} {}
pada suatu bundel vektor diferensiabel
\maabb {p} {E} {X
} {}
di atas suatu manifold diferensiabel $X$.
}

}
{} {}

\inputaufgabegibtloesung
{3}
{

Rumuskan teorema-teorema berikut.
\aufzaehlungdrei{\stichwort {Teorema isometri untuk transpor paralel pada suatu hiperpermukaan} {}

\mavergleichskette
{\vergleichskette
{  Y
}
{ \subseteq }{  \R^n
}
{  }{
}
{    }{
}
{   }{
}
}
{}{}{.}}{\stichwort {Theorema egregium} {.}}{Teorema tentang partisi kesatuan.}

}
{} {}

\inputaufgabegibtloesung
{3}
{

Misalkan
\maabbdisp {f,g} {\R} { \R^n
} {}
\definitionsverweis {kurva-kurva diferensiabel}{}{.}
Hitung
\definitionsverweis {turunan}{}{}
dari fungsi
\maabbeledisp {} {\R} {\R
} { t } { \left\langle f(t) , g(t)  \right\rangle
} {.}
Rumuskan hasilnya dengan menggunakan hasil kali dalam.

}
{} {}

\inputaufgabe
{0}
{

}
{} {}

\inputaufgabegibtloesung
{2}
{

Misalkan

\mavergleichskette
{\vergleichskette
{ Y
}
{ \subseteq }{ W
}
{ \subseteq }{ \R^n
}
{    }{
}
{   }{
}
}
{}{}{,}
$W$
\definitionsverweis {terbuka}{}{,}
suatu
\definitionsverweis {hiperpermukaan diferensiabel}{}{},
dan
\maabb {\gamma} {I =[a,b]} {Y
} {}
suatu
\definitionsverweis {kurva geodesik}{}{.}
Buktikan bahwa
\definitionsverweis {transpor paralel}{}{}
sepanjang $\gamma$ memetakan vektor tangensial

\mavergleichskette
{\vergleichskette
{ \gamma'(a)
}
{ \in }{ T_{\gamma(a)}Y
}
{  }{
}
{    }{
}
{   }{
}
}
{}{}{}
ke vektor tangensial

\mavergleichskette
{\vergleichskette
{ \gamma'(b)
}
{ \in }{ T_{\gamma(b)}Y
}
{  }{
}
{    }{
}
{   }{
}
}
{}{}{.}

}
{} {}

\inputaufgabe
{0}
{

}
{} {}

\inputaufgabegibtloesung
{3}
{

Apakah pemetaan
\maabbeledisp {\varphi} {S^1} {S^1
} {(x,y)} {( \betrag {  x  } ,y)
} {,}
\definitionsverweis {terdiferensialkan}{}{?}

}
{} {}

\inputaufgabegibtloesung
{7}
{

Misalkan $M$ suatu manifold diferensiabel kompak dengan bentuk volume positif kontinu $\omega$. Buktikan bahwa

\mavergleichskettedisp
{\vergleichskette
{ \int_M \omega
}
{ <} { \infty
}
{ } {
}
{ } {
}
{ } {
}
}
{}{}{}.

}
{} {}

\inputaufgabe
{0}
{

}
{} {}

\inputaufgabegibtloesung
{7 (1+2+4)}
{

\aufzaehlungdreiabc{ Berikan suatu sifat topologis yang menunjukkan bahwa
\definitionsverweis {interval satuan terbuka}{}{}
dan
\definitionsverweis {lingkaran}{}{}
$S^1$ tidak
\definitionsverweis {homeomorfik}{}{.}
}{Buktikan bahwa setiap
$1$-\definitionsverweis {bentuk diferensial}{}{}
kontinu pada
\mathl{]0,1[}{}
bersifat
\definitionsverweis {eksak}{}{.}
}{ Berikan suatu
$1$-bentuk konkret pada $S^1$ yang
\definitionsverweis {tertutup}{}{}
tetapi tidak eksak.
}

}
{} {}

\inputaufgabegibtloesung
{4}
{

Misalkan

\mavergleichskette
{\vergleichskette
{ U
}
{ \subseteq }{ \R^n
}
{  }{
}
{    }{
}
{   }{
}
}
{}{}{}
\definitionsverweis {terbuka}{}{}
dan
\maabb {f} { U } { \R
} {}
suatu
\definitionsverweis {fungsi yang terdiferensialkan secara kontinu}{}{}
dua kali. Buktikan bahwa

\mavergleichskettedisp
{\vergleichskette
{ d(df)
}
{ =} { 0
}
{ } {
}
{ } {
}
{ } {
}
}
{}{}{,}
dengan $d$ menyatakan
\definitionsverweis {turunan eksterior}{}{.}

}
{} {}

\inputaufgabegibtloesung
{3}
{

Deskripsikan
\definitionsverweis {pemetaan}{}{}
\maabbeledisp {\psi} { \Complex \setminus \{1\} } { \Complex
} {w} { { \frac{   (w+1)  { \mathrm i}  }{ 1-w  } }
} {}
dalam koordinat riil.

}
{} {}

\inputaufgabegibtloesung
{3}
{

Misalkan

\mavergleichskette
{\vergleichskette
{ H
}
{ \subseteq }{ \R^n
}
{  }{
}
{    }{
}
{   }{
}
}
{}{}{}
suatu
\definitionsverweis {setengah ruang Euklides}{}{}
dan

\mavergleichskette
{\vergleichskette
{ P
}
{ \in }{ H
}
{  }{
}
{    }{
}
{   }{
}
}
{}{}{.}
Andaikan terdapat suatu himpunan yang terbuka di $\R^n$, yaitu $V$, dengan

\mavergleichskette
{\vergleichskette
{ P
}
{ \in }{ V
}
{ \subseteq }{ H
}
{    }{
}
{   }{
}
}
{}{}{.}
Buktikan bahwa
\mathl{P}{} bukan titik batas dari $H$.

}
{} {}

\inputaufgabegibtloesung
{7}
{

Misalkan $M$ suatu
\definitionsverweis {manifold diferensiabel}{}{}
berdimensi $n$ dan

\mavergleichskette
{\vergleichskette
{ P
}
{ \in }{ M
}
{  }{
}
{    }{
}
{   }{
}
}
{}{}{}
suatu titik. Buktikan bahwa terdapat suatu
\definitionsverweis {manifold dengan batas}{}{}
$N$ yang batasnya $\partial N$ difeomorfik dengan permukaan bola berdimensi $(n-1)$, yaitu $S^{n-1}$, sedemikian sehingga
\mathkor {} {M \setminus \{ P \}} {dan} {N \setminus \partial N} {}
saling difeomorfik.

}
{} {}

\inputaufgabegibtloesung
{2}
{

Berikan suatu
\definitionsverweis {penghabisan kompak}{}{}
untuk $\R^d$.

}
{} {}

\inputaufgabegibtloesung
{9}
{

Buktikan teorema titik tetap Brouwer.

}
{} {}

\inputaufgabegibtloesung
{7 (1+2+4)}
{

Pada
\definitionsverweis {bundel vektor trivial}{}{}

\mathl{\R \times \R^2}{} di atas $\R$, kita meninjau
\definitionsverweis {koneksi trivial}{}{}
$\nabla$.
\aufzaehlungdrei{Buktikan bahwa

\mavergleichskette
{\vergleichskette
{ f_1(t)
}
{ = }{ \left(  1+ t^2  , \,  t^3                   \right)
}
{  }{
}
{    }{
}
{   }{
}
}
{}{}{}
dan

\mavergleichskette
{\vergleichskette
{ f_2(t)
}
{ = }{ \left(  t  , \,   1+t^2                   \right)
}
{  }{
}
{    }{
}
{   }{
}
}
{}{}{}
merupakan
\definitionsverweis {penampang basis}{}{}
dari bundel vektor tersebut.
}{Nyatakan penampang basis standar $e_1,e_2$ sebagai kombinasi linear dari penampang basis $f_1,f_2$.
}{Tentukan
\definitionsverweis {simbol-simbol Christoffel}{}{}
dari $\nabla$ terhadap penampang-penampang basis tersebut.
}

}
{} {}

 [[Kategorie:Latexseite]]
